\chapter*{Introduction générale}
\addcontentsline{toc}{chapter}{Introduction générale}
\markboth{Introduction générale}{Introduction générale}
\label{chap:introduction}

Chez Talan Tunisie, le processus de recrutement repose encore largement sur des
étapes manuelles : lecture de CV un par un, coordination d'entretiens par email,
évaluation des candidats sur des grilles papier. Face à un flux de candidatures
en constante augmentation, les recruteurs consacrent une part significative de
leur temps à des tâches répétitives, au détriment de l'analyse qualitative des
profils.

\medskip
Partant de ce constat, ce projet de fin d'études propose de concevoir et de
développer le \textbf{HR Talent Portal}, un portail de recrutement intelligent
construit sur \textbf{Salesforce Experience Cloud}. La particularité de cette
solution réside dans l'intégration de trois modèles d'intelligence artificielle
au c\oe{}ur du processus : \textbf{Groq LLaMA 3.3} pour le scoring automatique des
CV et l'analyse des entretiens d'architecture, et \textbf{Groq Whisper} pour la
transcription vocale en temps réel des visioconférences \textbf{Jitsi Meet}.

\medskip
Le portail couvre l'intégralité du cycle de recrutement, depuis le dépôt de
candidature jusqu'à la décision d'embauche, en passant par cinq étapes
automatisées : le screening IA du CV, un test de logique gamifié, un entretien
technique avec évaluation par IA, un entretien d'architecture avec
enregistrement audio et analyse automatique, puis la validation managériale.
Chaque acteur --- candidat, recruteur, manager et responsable RH --- dispose
de son propre tableau de bord avec une visibilité en temps réel sur l'avancement
du pipeline.

\medskip
Ce projet a été réalisé au sein du \textbf{centre d'excellence Salesforce de
Talan Tunisie}, en adoptant une méthodologie \textbf{Scrum} découpée en quatre
sprints. Le rapport s'organise en six chapitres :

\medskip
Le \textbf{premier chapitre} présente Talan Tunisie, l'écosystème Salesforce et
la problématique métier à l'origine du projet.

\medskip
Le \textbf{deuxième chapitre} détaille la planification : analyse des besoins
fonctionnels, modélisation UML et architecture technique retenue.

\medskip
Les quatre chapitres suivants correspondent aux sprints de réalisation :
\begin{itemize}
    \item \textbf{Sprint 1 :} Mise en place du portail Experience Cloud, gestion
    de l'identité candidat (inscription, authentification, reset mot de passe)
    et publication des offres d'emploi.
    \item \textbf{Sprint 2 :} Développement du scoring automatique des CV par
    Groq LLaMA, conversion des leads en candidatures et création automatique
    des comptes utilisateurs.
    \item \textbf{Sprint 3 :} Implémentation du test gaming, de l'entretien
    technique avec évaluation IA, et de l'entretien architecture avec
    visioconférence Jitsi, enregistrement audio, transcription Whisper et
    analyse LLaMA.
    \item \textbf{Sprint 4 :} Conception des tableaux de bord analytiques
    (recruteur, manager, RH), intégration du chatbot NLP via Agentforce,
    génération de rapports PDF, historique des candidatures et
    déclenchement du processus d'onboarding.
\end{itemize}
